\documentclass[11pt,a4paper]{article}
\usepackage[T1]{fontenc}
\usepackage[utf8]{inputenc}
\usepackage[provide=*,french]{babel}
\usepackage[a4paper,margin=2.3cm]{geometry}
\usepackage{hyperref}
\usepackage{enumitem}
\usepackage{booktabs,longtable,array}
\usepackage{xcolor}
\usepackage{fancyvrb}
\hypersetup{colorlinks=true, linkcolor=blue, urlcolor=blue, pdftitle={myco\_apps\_releases\_utilisateurs --- Guide rapide}}
\providecommand{\tightlist}{\setlength{\itemsep}{0pt}\setlength{\parskip}{0pt}}
\DefineVerbatimEnvironment{Highlighting}{Verbatim}{commandchars=\\\{\}}
\newenvironment{Shaded}{\begin{quote}}{\end{quote}}
\newcommand{\NormalTok}[1]{#1}
\newcommand{\ExtensionTok}[1]{#1}
\setlength{\emergencystretch}{3em}
\title{myco\_apps\_releases\_utilisateurs --- Guide rapide}
\author{Boite Eddy}
\date{\today}
\begin{document}
\maketitle
\section{myco\_apps\_releases\_utilisateurs --- Guide rapide}

\subsection{Objectif}

Cette archive permet de lancer localement deux analyses R, sans ouvrir RStudio :
\begin{itemize}
\item Inventaires fongiques --- Complétude \& Représentativité (ICR)
\item Évaluation potentiel fongique / intérêt patrimonial / gradient CHEGD
\end{itemize}

\subsection{Prérequis}

\begin{itemize}
\item R installé sur votre ordinateur : \url{https://cran.r-project.org/}
\item Connexion internet au premier lancement (installation automatique des packages)
\end{itemize}

\subsection{Lancement}

\subsubsection{Windows}
Double-cliquez : \texttt{Lancer\_Windows.bat}

\subsubsection{macOS}
Double-cliquez : \texttt{Lancer\_Mac.command}

Si macOS bloque le fichier, ouvrir un Terminal dans le dossier et exécuter :
\begin{verbatim}
xattr -dr com.apple.quarantine .
chmod +x Lancer_Mac.command
./Lancer_Mac.command
\end{verbatim}

\subsubsection{Linux}
Dans un terminal :
\begin{verbatim}
./Lancer_Linux.sh
\end{verbatim}

\subsection{Utilisation}

\begin{enumerate}
\item Choisir l'application à lancer.
\item Utiliser un fichier d'exemple ou importer votre propre fichier.
\item Cliquer sur \textbf{Lancer l'analyse}.
\item Consulter les sorties dans l'onglet \textbf{Résultats} ou dans \texttt{applications/results/}.
\end{enumerate}

\subsection{Colonnes attendues}

\subsubsection{ICR}
Colonnes obligatoires : \texttt{site}, \texttt{date}, \texttt{visite\_id}, \texttt{espece}.

Colonne optionnelle : \texttt{placette}.

\subsubsection{CHEGD}
Colonnes attendues :
\begin{itemize}
\item \texttt{Espèces}, \texttt{Famille}, \texttt{Date} ;
\item \texttt{Nombre d'espèce}, \texttt{Site}, \texttt{Fiabilité détermination}.
\end{itemize}

\subsection{Où trouver les résultats et les logs ?}

\begin{itemize}
\item Résultats ICR : \texttt{applications/results/ICR/}
\item Résultats CHEGD : \texttt{applications/results/EPFIP\_CHEGD/}
\item Logs d'exécution : \texttt{applications/logs/}
\end{itemize}

En cas de problème, transmettre le dernier fichier \texttt{run\_...log}.

\subsection{Fin de session}

Dans l'application, cliquer sur \textbf{Quitter l'application}. Fermer ensuite l'onglet navigateur et la fenêtre Terminal si nécessaire.

\subsection{Contenu recommandé du ZIP}

Distribuer uniquement le dossier :
\begin{verbatim}
myco_apps_releases_utilisateurs/
\end{verbatim}

Ce dossier est autonome pour le déploiement utilisateur.

\end{document}
